\scanpage{25}
\mspage{24}
% @ opening
Mais, \unsure{on|\deleted{on}} en arrive aux réciproques

% @ upper left, oblique
% slanted
\margin{Attention, Ces réciproques valent dans les exemples-types \unsure{unistrates|uni-strates} (théorie topologique, ou modérée) mais pas pour les réseaux p.ex.}

% @ Cont 5
\uline{Cont} 5 Soient $X,Y\in M$ deux multistrates, et considérons $\deleted{\unsure{\mathcal{L}_X|\widehat{\mathcal{L}}_X}}\widehat{\mathcal{L}}_X,\widehat{\mathcal{L}}_Y\subset\mathcal P(\mathcal{L})$ les ens de parties de $\mathcal{L}$ correspondantes. Ils permettent, grâce aux scholies de récupérer les strates ``ouvertes''

% @ a), first condition
a) Pour que $X\ll Y$, il f. et s. que \deleted{tout $\forall A\ \mathcal{L}\ \forall A\in \mathcal{L}_X$} toute strate ouverte de $X$ soit contenue dans une strate ouverte de $Y$

% @ left, beside (a) and (b)
% slanted
\margin{Il en résulte l'énoncé correspondant, pour deux figures $F,F'$, \unsure{déterminé|déterminés|déterminées} \unsure{par} $\widehat{\mathcal{L}}_F,\widehat{\mathcal{L}}_{F'}$.}

% @ b), second condition
b) Pour que $X$ et $Y$ soient compatibles, il f. et s. que pour toute strate ouverte de $X$ et toute strate ouverte de $Y$, \deleted{on ait} celles-ci soient disjointes ou identiques.

% @ Cor., two lines
\uline{Cor.}
\[
\begin{aligned}
\text{\deleted{b)} a)}\quad &\widehat{\mathcal{L}}_X\ne\widehat{\mathcal{L}}_Y\ \deleted{\unsure{\text{donc}|\Longrightarrow}}\Longleftrightarrow\ X=Y\\
\text{b)}\quad &\widehat{\mathcal{L}}_X\subset\widehat{\mathcal{L}}_Y\ \deleted{\unsure{\widehat{\mathcal{L}}_X\,\widehat{\mathcal{L}}_X|\widehat{\mathcal{L}}_X}}\ \text{est fermé dans }\widehat{\mathcal{L}}_Y\ \Longleftrightarrow\ X\leq Y
\end{aligned}
\]

% @ Scholie, bottom
\uline{Scholie} On récupère $M$, avec $\ll$, $\mathcal R$ (compatibilité), en termes de $\mathcal{L}$, et de l'ens \inserted{$\widehat M$} des $\widehat{\mathcal{L}}_X$
\[
\widehat M=\operatorname{Im}\bigl(\underbrace{M\xrightarrow{X\mapsto\widehat{\mathcal{L}}_X}\mathcal P(\mathcal P}_{\text{application \uline{injective}}}(\mathcal{L}))\ \in\mathcal P(\mathcal P(\mathcal P(\mathcal{L})))
\]
